\documentclass[12pt]{article}
\usepackage[utf8]{inputenc}
\usepackage{amsmath, amssymb}
\usepackage[spanish]{babel}
\usepackage{geometry}
\geometry{a4paper, margin=1.5cm}
\usepackage{fancyhdr}
\usepackage{eso-pic}
\usepackage{graphicx}
\usepackage{booktabs}   % Para tablas profesionales
\usepackage{longtable}  % Para tablas que pueden ocupar varias páginas

% Configuración de página y encabezados/pies
\pagestyle{fancy}
\fancyhf{} % limpia todo
\fancyfoot[R]{Christian Cristobal Carbajal García} % autor abajo a la derecha
\renewcommand{\headrulewidth}{1pt} % línea en el encabezado
\renewcommand{\footrulewidth}{1pt} % línea en el pie

% Logo opcional en esquina superior izquierda (solo primera página)
\makeatletter
\AddToShipoutPictureFG*{%
  \AtPageUpperLeft{%
    \put(20,-60){%
      \IfFileExists{/home/christian/Documentos/Reportes_Aut/imagenes/logo_pdf.png}{%
        
\includegraphics[width=2.5cm]{/home/christian/Documentos/Reportes_Aut/imagenes/logo_pdf.png}%
      }{%
        % Si no existe el archivo, no hace nada (sin error)
      }%
    }%
  }%
}
\makeatother

% =========================
% IMAGEN SUPERIOR DERECHA (CORREGIDO)
% =========================
\AddToShipoutPictureFG*{%
  \AtPageUpperLeft{%
    \put(500,-85){% <-- ajusta este valor si lo quieres mover
      \IfFileExists{/home/christian/Documentos/Reportes_Aut/imagenes/marca_de_agua.png}{%
        \includegraphics[width=4.0cm]{/home/christian/Documentos/Reportes_Aut/imagenes/marca_de_agua.png}%
      }{}%
    }%
  }%
}



\title{\textbf{Tabla de Constantes Físicas}}
\author{}
\date{}

\begin{document}

\maketitle

\noindent
Esta tabla recopila las constantes físicas más utilizadas en diversos campos de la física. Los valores están expresados en unidades del Sistema Internacional (SI), salvo indicación contraria.

\begin{longtable}{p{0.35\textwidth} p{0.2\textwidth} p{0.35\textwidth}}
\caption{Constantes físicas fundamentales y derivadas}\\
\toprule
\textbf{Constante} & \textbf{Símbolo} & \textbf{Valor (unidades)} \\
\midrule
\endfirsthead

\multicolumn{3}{c}%
{{\bfseries \tablename\ \thetable{} -- Continuación}} \\
\toprule
\textbf{Constante} & \textbf{Símbolo} & \textbf{Valor (unidades)} \\
\midrule
\endhead

\bottomrule
\multicolumn{3}{r}{{Continúa en la página siguiente}} \\
\endfoot

\bottomrule
\endlastfoot

% Masa de partículas alfa
Masa de la partícula $\alpha$ & $m_\alpha$ & $6.6446573357 \times 10^{-27}\ \text{kg}$ \\
% Constante de masa atómica (unidad de masa atómica unificada)
Constante de masa atómica & $u$ & $1.66053906660 \times 10^{-27}\ \text{kg}$ \\
% Unidad atómica de energía (Hartree)
Unidad atómica de energía & $E_h$ & $4.3597447222071 \times 10^{-18}\ \text{J}$ \\
% Unidad atómica de fuerza
Unidad atómica de fuerza & $F_{\text{au}}$ & $8.2387234983 \times 10^{-8}\ \text{N}$ \\
% Unidad atómica de longitud (radio de Bohr)
Unidad atómica de longitud & $a_0$ & $5.29177210903 \times 10^{-11}\ \text{m}$ \\
% Unidad atómica de masa (masa del electrón)
Unidad atómica de masa & $m_e$ & $9.1093837015 \times 10^{-31}\ \text{kg}$ \\
% Constante de Avogadro
Constante de Avogadro & $N_A$ & $6.02214076 \times 10^{23}\ \text{mol}^{-1}$ \\
% Radio de Bohr
Radio de Bohr & $a_0$ & $5.29177210903 \times 10^{-11}\ \text{m}$ \\
% Constante de Boltzmann
Constante de Boltzmann & $k_B$ & $1.380649 \times 10^{-23}\ \text{J/K}$ \\
% Radio clásico del electrón
Radio clásico del electrón & $r_e$ & $2.8179403262 \times 10^{-15}\ \text{m}$ \\
% Carga elemental
Carga elemental & $e$ & $1.602176634 \times 10^{-19}\ \text{C}$ \\
% Masa de la Tierra
Masa de la Tierra & $M_\oplus$ & $5.9722 \times 10^{24}\ \text{kg}$ \\
% Radio ecuatorial de la Tierra
Radio ecuatorial de la Tierra & $R_\oplus$ & $6.3781 \times 10^{6}\ \text{m}$ \\
% Velocidad orbital media de la Tierra
Velocidad orbital media de la Tierra & $v_\oplus$ & $2.978 \times 10^{4}\ \text{m/s}$ \\
% Velocidad angular de rotación de la Tierra
Velocidad angular de rotación de la Tierra & $\omega_\oplus$ & $7.292115 \times 10^{-5}\ \text{rad/s}$ \\
% Constante de Faraday
Constante de Faraday & $F$ & $9.648533212 \times 10^{4}\ \text{C/mol}$ \\
% Constante de los gases ideales
Constante de los gases ideales & $R$ & $8.314462618 \ \text{J/(mol·K)}$ \\
% Constante de masa molar
Constante de masa molar & $M_u$ & $1 \times 10^{-3}\ \text{kg/mol}$ \\
% Volumen molar del gas ideal (0°C, 1 atm)
Volumen molar del gas ideal & $V_m$ & $2.241396 \times 10^{-2}\ \text{m}^3/\text{mol}$ \\
% Masa del neutrón
Masa del neutrón & $m_n$ & $1.674927498 \times 10^{-27}\ \text{kg}$ \\
% Constante de Newton de gravitación
Constante de gravitación universal & $G$ & $6.67430 \times 10^{-11}\ \text{m}^3\ \text{kg}^{-1}\ \text{s}^{-2}$ \\
% Constante de Planck
Constante de Planck & $h$ & $6.62607015 \times 10^{-34}\ \text{J·s}$ \\
% Masa de Planck
Masa de Planck & $m_P$ & $2.176434 \times 10^{-8}\ \text{kg}$ \\
% Masa del protón
Masa del protón & $m_p$ & $1.67262192369 \times 10^{-27}\ \text{kg}$ \\
% Relación masa protón-electrón
Relación masa protón-electrón & $m_p/m_e$ & $1836.15267343$ \\
% Velocidad de la luz en el vacío
Velocidad de la luz en el vacío & $c$ & $2.99792458 \times 10^{8}\ \text{m/s}$ \\
% Aceleración normal de la gravedad
Aceleración normal de la gravedad & $g_n$ & $9.80665\ \text{m/s}^2$ \\
% Atmósfera estándar
Atmósfera estándar & $1\ \text{atm}$ & $1.01325 \times 10^{5}\ \text{Pa}$ \\
% Referencia estándar de intensidad de sonido
Intensidad de referencia (sonido) & $I_0$ & $1 \times 10^{-12}\ \text{W/m}^2$ \\
% Constante de Stefan-Boltzmann
Constante de Stefan-Boltzmann & $\sigma$ & $5.670374419 \times 10^{-8}\ \text{W·m}^{-2}\ \text{K}^{-4}$ \\
\end{longtable}

\vspace{1cm}
\noindent
\textbf{Nota:} Los valores corresponden a las recomendaciones CODATA 2018 y ajustes posteriores (2022). Las constantes definidas exactamente (como $c$, $h$, $e$, $k_B$, $N_A$) tienen valores fijos desde la redefinición del SI en 2019.

\end{document}